\textbf{Phần 1. Trắc nghiệm lựa chọn}
\setcounterpageref{ex}{0}

\begin{ex} %[2X14N74]
Trong một túi chỉ có các viên bi màu xanh. Lấy ra 1 viên bi, sự kiện "Lấy được viên bi xanh" là
\choice
{Có thể}
{Không thể}
{\True Chắc chắn}
{Chắc chắn không}
\loigiai{
Chỉ có bi xanh nên chắc chắn lấy được bi xanh.\\
\textbf{Đáp án của câu hỏi là C}
}
\end{ex}

\begin{ex} %[2X14N74]
Trong một túi chỉ có các viên bi màu xanh. Lấy ra 1 viên bi, sự kiện "Lấy được viên bi đỏ" là
\choice
{Chắc chắn}
{Có thể}
{Có thể hoặc không}
{\True Không thể}
\loigiai{
Không có bi đỏ nên không thể lấy được bi đỏ.\\
\textbf{Đáp án của câu hỏi là D}
}
\end{ex}

\begin{ex} %[2X14H74]
Trong túi có cả kẹo dâu và kẹo cam. Lấy ra 1 cái kẹo, sự kiện "Lấy được kẹo dâu" là
\choice
{\True Có thể}
{Chắc chắn}
{Không thể}
{Chắc chắn không}
\loigiai{
Có thể lấy được kẹo dâu.\\
\textbf{Đáp án của câu hỏi là A}
}
\end{ex}

\begin{ex} %[2X14H74]
Kiểm đếm các loại hoa trong lọ: 6 hoa hồng, 4 hoa cúc, 5 hoa đồng tiền. Loại hoa có nhiều nhất là
\choice
{Hoa cúc}
{Hoa đồng tiền}
{Bằng nhau}
{\True Hoa hồng ($6$ bông)}
\loigiai{
Hoa hồng có 6 bông (nhiều nhất).\\
\textbf{Đáp án của câu hỏi là D}
}
\end{ex}

\begin{ex} %[2X14H74]
Tổng số hoa có trong lọ là
\choice
{$14$ bông}
{\True $15$ bông}
{$16$ bông}
{$17$ bông}
\loigiai{
$6 + 4 + 5 = 15$ bông hoa.\\
\textbf{Đáp án của câu hỏi là B}
}
\end{ex}

\begin{ex} %[2X14H74]
Số hoa hồng nhiều hơn số hoa cúc là
\choice
{\True $2$ bông}
{$1$ bông}
{$3$ bông}
{$4$ bông}
\loigiai{
$6 - 4 = 2$ bông.\\
\textbf{Đáp án của câu hỏi là A}
}
\end{ex}

\begin{ex} %[2X14V74]
Biểu đồ tranh ghi số sách của mỗi bạn (mỗi biểu tượng $= 2$ cuốn sách): An có 3 biểu tượng, Bình có 5 biểu tượng. Số sách của Bình là
\choice
{$8$ cuốn}
{$12$ cuốn}
{\True $10$ cuốn}
{$15$ cuốn}
\loigiai{
$2 \times 5 = 10$ cuốn.\\
\textbf{Đáp án của câu hỏi là C}
}
\end{ex}

\begin{ex} %[2X14V74]
Bình có nhiều hơn An số cuốn sách là
\choice
{$2$ cuốn}
{\True $4$ cuốn}
{$6$ cuốn}
{$8$ cuốn}
\loigiai{
An có: $2 \times 3 = 6$ cuốn.\\
$10 - 6 = 4$ cuốn.\\
\textbf{Đáp án của câu hỏi là B}
}
\end{ex}

\begin{ex} %[2X14V74]
Gieo một con xúc xắc 6 mặt. Sự kiện "Số chấm xuất hiện lớn hơn 0 và nhỏ hơn 7" là
\choice
{Có thể}
{Không thể}
{Tùy lần gieo}
{\True Chắc chắn}
\loigiai{
Xúc xắc chỉ có từ 1 đến 6 chấm nên chắc chắn xảy ra.\\
\textbf{Đáp án của câu hỏi là D}
}
\end{ex}

\begin{ex} %[2X14V74]
Thực hành quay chong chóng 20 lần: kim chỉ vào màu đỏ 12 lần, màu xanh 8 lần. Số lần kim chỉ vào màu đỏ nhiều hơn màu xanh là
\choice
{\True $4$ lần}
{$2$ lần}
{$6$ lần}
{$8$ lần}
\loigiai{
$12 - 8 = 4$ lần.\\
\textbf{Đáp án của câu hỏi là A}
}
\end{ex}

\begin{ex} %[2X14V74]
Khảo sát trái cây yêu thích của 35 học sinh: 15 bạn thích cam, 12 bạn thích xoài, còn lại thích táo. Số bạn thích táo là
\choice
{$7$ bạn}
{\True $8$ bạn}
{$9$ bạn}
{$10$ bạn}
\loigiai{
$35 - (15 + 12) = 35 - 27 = 8$ bạn.\\
\textbf{Đáp án của câu hỏi là B}
}
\end{ex}

\begin{ex} %[2X14C74]
Trong hộp có 4 viên bi đỏ và 4 viên bi xanh. Lấy ra ngẫu nhiên 5 viên bi. Khẳng định nào sau đây là CHẮC CHẮN?
\choice
{Lấy được cả 5 viên bi đỏ}
{Lấy được cả 5 viên bi xanh}
{\True Trong 5 viên lấy ra có ít nhất 1 viên bi đỏ và 1 viên bi xanh}
{Lấy được 3 viên bi vàng}
\loigiai{
Vì mỗi màu chỉ có 4 viên nên lấy 5 viên chắc chắn có cả hai màu.\\
\textbf{Đáp án của câu hỏi là C}
}
\end{ex}

\textbf{Phần 2. Trắc nghiệm đúng sai}
\setcounterpageref{ex}{0}

\begin{ex} %[2X14H74]
Trong một hộp có 6 bút chì đen và 2 bút chì đỏ. Xét tính đúng sai khi lấy ra 1 chiếc bút:
\choiceTF
{\True Có thể lấy được bút chì đen.}
{\True Có thể lấy được bút chì đỏ.}
{Chắc chắn lấy được bút chì đen.}
{\True Không thể lấy được bút chì màu xanh (vì không có bút xanh).}
\loigiai{
\begin{itemize}
    \item Có thể lấy bút đen: Đúng.
    \item Có thể lấy bút đỏ: Đúng.
    \item Không chắc chắn đen vì có thể ra đỏ. Ý 3 Sai, Ý 4 Đúng.
\end{itemize}
\begin{itemize}
    \item a) \textbf{Đúng}.
    \item b) \textbf{Đúng}.
    \item c) \textbf{Sai}.
    \item d) \textbf{Đúng}.
\end{itemize}
}
\end{ex}

\begin{ex} %[2X14V74]
Biểu đồ tranh số con gà của 2 bạn (mỗi biểu tượng $= 5$ con gà): Nam có 4 biểu tượng, Hải có 3 biểu tượng. Xét tính đúng sai:
\choiceTF
{\True Nam có 20 con gà ($5 \times 4 = 20$).}
{\True Hải có 15 con gà ($5 \times 3 = 15$).}
{Cả hai bạn có tất cả 40 con gà.}
{\True Cả hai bạn có tất cả 35 con gà ($20 + 15 = 35$).}
\loigiai{
\begin{itemize}
    \item Nam: 20 con: Đúng.
    \item Hải: 15 con: Đúng.
    \item Tổng: $20 + 15 = 35 \neq 40$. Ý 3 Sai, Ý 4 Đúng.
\end{itemize}
\begin{itemize}
    \item a) \textbf{Đúng}.
    \item b) \textbf{Đúng}.
    \item c) \textbf{Sai}.
    \item d) \textbf{Đúng}.
\end{itemize}
}
\end{ex}

\textbf{Phần 3. Trả lời ngắn}
\setcounterpageref{ex}{0}

\begin{ex} %[2X14H74]
Trong hộp có 5 viên bi xanh. Nhắm mắt lấy 1 viên bi, sự kiện "Lấy được bi xanh" là gì? (Viết: chắc chắn, có thể, hoặc không thể)
\shortans{chắc chắn}
\loigiai{
\textbf{Đáp số:} chắc chắn\\
Chắc chắn.
}
\end{ex}

\begin{ex} %[2X14H74]
Trong hộp có 5 viên bi xanh. Nhắm mắt lấy 1 viên bi, sự kiện "Lấy được bi vàng" là gì? (Viết: chắc chắn, có thể, hoặc không thể)
\shortans{không thể}
\loigiai{
\textbf{Đáp số:} không thể\\
Không thể.
}
\end{ex}

\begin{ex} %[2X14V74]
Biểu đồ tranh: mỗi hình bông hoa biểu diễn cho 2 bông hoa thật. Có 6 hình bông hoa, hỏi có bao nhiêu bông hoa thật?
\shortans{12}
\loigiai{
\textbf{Đáp số:} $12$\\
$2 \times 6 = 12$ bông.
}
\end{ex}

\begin{ex} %[2X14V74]
Khảo sát 30 học sinh: 14 bạn thích đá bóng, 10 bạn thích bơi, còn lại thích cờ vua. Hỏi có bao nhiêu bạn thích cờ vua?
\shortans{6}
\loigiai{
\textbf{Đáp số:} $6$\\
$30 - (14 + 10) = 30 - 24 = 6$ bạn.
}
\end{ex}

\begin{ex} %[2X14V74]
Gieo đồng xu 15 lần thu được 9 lần mặt Sấp. Hỏi mặt Ngửa xuất hiện bao nhiêu lần?
\shortans{6}
\loigiai{
\textbf{Đáp số:} $6$\\
$15 - 9 = 6$ lần.
}
\end{ex}

\begin{ex} %[2X14C74]
Có 50 cây xanh cần biểu diễn lên biểu đồ tranh. Biết mỗi biểu tượng cây biểu diễn cho 5 cây thật. Hỏi cần vẽ bao nhiêu biểu tượng cây?
\shortans{10}
\loigiai{
\textbf{Đáp số:} $10$\\
$50 : 5 = 10$ biểu tượng.
}
\end{ex}
